\section{Introdução e Motivação}

\begin{frame}{Contextualização do Problema}

\begin{itemize}

    \item Modelos matemáticos baseados em
    \textbf{equações diferenciais ordinárias (EDOs)}
    são amplamente utilizados para descrever
    a evolução temporal de sistemas dinâmicos.

    \vspace{0.3cm}

    \item A capacidade preditiva desses modelos
    depende da determinação adequada de seus
    \textbf{parâmetros}.

    \vspace{0.3cm}

    \item A identificação desses parâmetros a partir
    de observações constitui um
    \textbf{problema inverso}, frequentemente
    formulado como um problema de otimização.

    \vspace{0.3cm}

    \item Sua resolução pode apresentar desafios,
    como:
    \begin{itemize}
        \item Sensibilidade ao ruído e às condições iniciais;
        \item Não linearidade e mínimos locais;
        \item Identificabilidade dos parâmetros;
        \item Custo computacional.
    \end{itemize}

\end{itemize}

\end{frame}

\begin{frame}{Objetivos do Trabalho}

\textbf{Objetivo Geral}

\begin{itemize}
    \item Investigar e comparar diferentes técnicas de
    \textbf{identificação de parâmetros} em problemas
    inversos, utilizando o modelo de Rosenzweig-MacArthur
    como estudo de caso.
\end{itemize}

\vspace{0.35cm}

\textbf{Objetivos Específicos}

\begin{itemize}

    \item Implementar um algoritmo genético para
    calibração de parâmetros;

    \vspace{0.15cm}

    \item Investigar abordagens baseadas em
    \textbf{redes neurais}, incluindo redes neurais
    inversas e PINNs;

    \vspace{0.15cm}

    \item Avaliar a capacidade de recuperação dos
    parâmetros a partir de dados sintéticos,
    com diferentes níveis de ruído;

    \vspace{0.15cm}

    \item Comparar as técnicas quanto à
    \textbf{precisão, robustez, sensibilidade}
    e \textbf{custo computacional}.

\end{itemize}

\end{frame}

\section{Definição do Problema Inverso}

\begin{frame}{Modelo de Rosenzweig-MacArthur}

O modelo de \textbf{Rosenzweig-MacArthur} descreve
a dinâmica entre presas e predadores, incorporando o
crescimento logístico das presas e a resposta funcional
de Holling tipo II.

\vspace{0.25cm}

\begin{columns}
    \begin{column}{0.5\textwidth}
        \begin{block}{Sistema de Equações Diferenciais}
        \begin{equation}
        \begin{cases}
        \dfrac{dN}{dt} = rN\left(1-\dfrac{N}{K}\right)
        - \dfrac{aNP}{1+ahN}, \\[8pt]
        \dfrac{dP}{dt} = e\dfrac{aNP}{1+ahN} - mP.
        \end{cases}
        \end{equation}
        \end{block}
    \end{column}
    \begin{column}{0.5\textwidth}
        Em que $N(t)$ e $P(t)$ representam, respectivamente,
        as populações de presas e predadores.

        \vspace{0.2cm}

        \begin{itemize}
            \item $r$: taxa de crescimento das presas;
            \item $K$: capacidade de suporte;
            \item $a$: taxa de ataque dos predadores;
            \item $h$: tempo de manipulação;
            \item $e$: eficiência de conversão;
            \item $m$: taxa de mortalidade dos predadores.
        \end{itemize}
    \end{column}
\end{columns}

\vspace{0.15cm}

O vetor de parâmetros a ser identificado é:
\[
\boldsymbol{\theta} =
(r,K,a,h,e,m).
\]

\end{frame}

\begin{frame}{Geração dos Dados Sintéticos}

    \begin{itemize}
        \item Definir um conjunto de parâmetros de referência:
        \[
        \boldsymbol{\theta}^{*}
        =
        (r^{*},K^{*},a^{*},h^{*},e^{*},m^{*}).
        \]

        \item Resolver numericamente o sistema de
        Rosenzweig-MacArthur, considerando condições iniciais
        conhecidas:
        \[
        N(0)=N_0,
        \qquad
        P(0)=P_0.
        \]

        \item Amostrar as soluções em instantes discretos
        $t_1,t_2,\ldots,t_m$ e adicionar ruído:
        \[
        \begin{aligned}
        N_i^{obs} &= N(t_i;\boldsymbol{\theta}^{*})
                    + \varepsilon_{N,i},\\
        P_i^{obs} &= P(t_i;\boldsymbol{\theta}^{*})
                    + \varepsilon_{P,i}.
        \end{aligned}
        \]

    \end{itemize}
\end{frame}

\begin{frame}{Geração de Dados Sintéticos}
    \begin{figure}
        \includegraphics[width=0.8\textwidth]{imgs/dados.png}
    \end{figure}
\end{frame}

\begin{frame}{Geração de Dados Sintéticos}
    
    Considerar diferentes cenários experimentais:
    \begin{itemize}
        \item Dados sem ruído e com diferentes níveis de ruído;
        \item Diferentes quantidades de observações temporais;
        \item Diferentes conjuntos de parâmetros de referência.
    \end{itemize}

\vspace{0.15cm}

\begin{block}{Finalidade}
    Construir problemas inversos controlados, permitindo
    avaliar a capacidade de cada método em recuperar
    os parâmetros conhecidos.
\end{block}
\end{frame}

\begin{frame}{Formulação do Problema Inverso}

\begin{columns}
    \begin{column}{0.65\textwidth}
        \begin{itemize}

            \item Considere as observações disponíveis em
            $m$ instantes de tempo:
            \[
            \mathcal{D}
            =
            \left\{
            (t_i,N_i^{obs},P_i^{obs})
            \right\}_{i=1}^{m}.
            \]

            \item O objetivo é identificar o vetor de parâmetros:
            \[
            \boldsymbol{\theta}
            =
            (r,K,a,h,e,m),
            \]
            cujas soluções numéricas melhor reproduzam
            as observações.

        \end{itemize}
        
    \end{column}
    \begin{column}{0.3\textwidth}
        \begin{itemize}
        \item $\Omega$: conjunto admissível de parâmetros positivos;
        \item $w_N,w_P$: pesos para balancear os resíduos;
        \item $N(t;\boldsymbol{\theta}),P(t;\boldsymbol{\theta})$:
        soluções numéricas do modelo.
        \end{itemize}
    \end{column}
\end{columns}

\begin{block}{Problema de Otimização}
        \[
        \boxed{
        \boldsymbol{\theta}^{opt}
        =
        \underset{\boldsymbol{\theta}\in\Omega}{\arg\min}
        \; J(\boldsymbol{\theta})
        }
        \]

        \vspace{-0.2cm}

        \[
        J(\boldsymbol{\theta})
        =
        \frac{1}{m}
        \sum_{i=1}^{m}
        \left[
        w_N\left(N(t_i;\boldsymbol{\theta})-N_i^{obs}\right)^2
        +
        w_P\left(P(t_i;\boldsymbol{\theta})-P_i^{obs}\right)^2
        \right].
        \]
\end{block}
\end{frame}

\section{Métodos de Identificação de Parâmetros}

\begin{frame}{Algoritmo Genético}

\small O AG explora populações de parâmetros candidatos sem exigir o
cálculo de derivadas da função objetivo.

\vspace{0.2cm}

\begin{center}
\begin{tikzpicture}[
    >=stealth,
    etapa/.style={
        draw=LNCCmainblue,
        fill=LNCClighteraltblue,
        rounded corners=3pt,
        align=center,
        minimum width=2.25cm,
        minimum height=0.9cm,
        font=\scriptsize
    },
    decisao/.style={
        draw=LNCCdarkorange,
        fill=LNCClightyellow,
        diamond,
        aspect=1.7,
        align=center,
        inner sep=1pt,
        font=\scriptsize
    },
    resultado/.style={
        draw=LNCCdarkgreen,
        fill=LNCClightgreen,
        rounded corners=3pt,
        align=center,
        minimum width=2.25cm,
        minimum height=0.75cm,
        font=\scriptsize
    },
    seta/.style={->, thick, draw=LNCCdarkaltblue}
]
    \node[etapa] (populacao) at (0,0) {População inicial\\$\boldsymbol{\theta}_i=(r,K,a,h,e,m)$};
    \node[etapa] (avaliacao) at (3.1,0) {Resolver as EDOs\\e calcular $J(\boldsymbol{\theta}_i)$};
    \node[etapa] (selecao) at (6.2,0) {Seleção dos\\melhores indivíduos};
    \node[etapa] (operadores) at (9.3,0) {Cruzamento e mutação\\geram nova população};
    \node[decisao] (parada) at (9.3,-1.65) {Critério de\\parada?};
    \node[resultado] (resultado) at (9.3,-2.85) {Melhor solução\\$\hat{\boldsymbol{\theta}}$};

    \draw[seta] (populacao) -- (avaliacao);
    \draw[seta] (avaliacao) -- (selecao);
    \draw[seta] (selecao) -- (operadores);
    \draw[seta] (operadores) -- (parada) node[midway, right, font=\scriptsize] {nova geração};
    \draw[seta] (parada) -- (resultado) node[midway, right, font=\scriptsize] {sim};
    \draw[seta] (parada.west) -| (avaliacao.south) node[pos=0.18, below, font=\scriptsize] {não};
\end{tikzpicture}
\end{center}

\end{frame}

\begin{frame}{Redes Neurais para Problemas Inversos}

\begin{figure}
    \includegraphics[width=0.7\textwidth]{imgs/NN.png}
\end{figure}
\end{frame}

\begin{frame}{Rede Neural Inversa: Dados e Características}
\scriptsize
\begin{columns}[T,onlytextwidth]
    \begin{column}{0.3\textwidth}
        \begin{block}{Dados utilizados}
        \begin{itemize}
            \item Banco supervisionado de pares
            $(\mathbf{y}^{(j)},\boldsymbol{\theta}^{(j)})$;
            \item Trajetórias geradas para diferentes parâmetros,
            condições iniciais e níveis de ruído;
            \item Separação em treino, validação e teste por
            vetores de parâmetros.
        \end{itemize}
        \end{block}
    \end{column}
    \begin{column}{0.3\textwidth}
        \begin{block}{Vantagens}
        \begin{itemize}
            \item Inferência rápida após o treinamento;
            \item Reutilizável para muitas trajetórias observadas;
            \item Não requer resolver as EDOs em cada nova
            estimação.
        \end{itemize}
        \end{block}
    \end{column}
    \begin{column}{0.32\textwidth}
        \begin{alertblock}{Limitações}
        \begin{itemize}
            \item Banco de treinamento grande e representativo;
            \item Custo inicial de gerar trajetórias rotuladas;
            \item Generalização limitada fora do domínio treinado.
        \end{itemize}
        \end{alertblock}
    \end{column}
\end{columns}

\vspace{0.05cm}
\begin{block}{Uso no teste}
Uma nova trajetória observada, associada a parâmetros desconhecidos,
é fornecida à rede para obter $\hat{\boldsymbol{\theta}}$.
\end{block}
\end{frame}

\begin{frame}{Physics-Informed Neural Networks (PINNs)}

\begin{figure}
    \includegraphics[width=0.7\textwidth]{imgs/PINN.png}
\end{figure}
\end{frame}

\begin{frame}{PINN Inversa: Dados e Características}
\scriptsize
\begin{columns}[T,onlytextwidth]
    \begin{column}{0.3\textwidth}
        \begin{block}{Dados utilizados}
        \begin{itemize}
            \item Uma única série observada de presas e predadores;
            \item Condições iniciais conhecidas ou observadas;
            \item Pontos de colocação no tempo para impor as EDOs.
        \end{itemize}
        \end{block}
    \end{column}
    \begin{column}{0.3\textwidth}
        \begin{block}{Vantagens}
        \begin{itemize}
            \item Não exige banco prévio de trajetórias rotuladas;
            \item Integra observações e a física do modelo;
            \item Reconstrói trajetórias contínuas no tempo.
        \end{itemize}
        \end{block}
    \end{column}
    \begin{column}{0.32\textwidth}
        \begin{alertblock}{Limitações}
        \begin{itemize}
            \item Novo treinamento para cada conjunto de observações;
            \item Custo elevado e sensibilidade aos pesos das perdas;
            \item Ruído e baixa identificabilidade dificultam a estimação.
        \end{itemize}
        \end{alertblock}
    \end{column}
\end{columns}

\vspace{0.05cm}
\begin{block}{Avaliação com dados sintéticos}
O parâmetro verdadeiro é usado apenas para medir o erro de
$\hat{\boldsymbol{\theta}}$, e não como rótulo durante o treinamento.
\end{block}
\end{frame}

\section{Metodologia Experimental}

\begin{frame}{Protocolo Experimental e Número de Otimizações}
\small
\begin{columns}[T,onlytextwidth]
    \begin{column}{0.42\textwidth}
        \begin{block}{Cenários de teste}
        \begin{itemize}
            \item Variar $N_\theta$ vetores de parâmetros,
            $N_\sigma$ níveis de ruído e $N_m$ quantidades
            de observações;
            \item Total de cenários:
            \[
            N_c=N_\theta N_\sigma N_m;
            \]
            \item Repetir cada cenário $R$ vezes, com sementes
            e inicializações distintas.
        \end{itemize}
        \end{block}
    \end{column}
    \begin{column}{0.55\textwidth}
        \begin{block}{Execuções por método}
        \centering
        \scriptsize
        \begin{tabular}{@{}p{0.24\linewidth}p{0.48\linewidth}c@{}}
        \toprule
        Método & Por caso & Total \\
        \midrule
        AG & Uma calibração evolutiva & $N_cR$ \\
        Rede neural inversa & Treino único; inferência nos testes & $N_{NN}$ \\
        PINN & Uma PINN por série observada & $N_cR$ \\
        \bottomrule
        \end{tabular}
        \end{block}
    \end{column}
\end{columns}

\vspace{0.1cm}
\begin{alertblock}{Contagem total}
\[
N_{\mathrm{otim}} = 2N_cR + N_{NN},
\]
em que $N_{NN}$ é o número de treinamentos da rede inversa
(configurações de arquitetura e sementes). Para uma única configuração
e semente, $N_{NN}=1$.
\end{alertblock}
\end{frame}

\section{Estratégia de Avaliação}

\begin{frame}{Métricas de Avaliação}

\small

\begin{itemize}

    \item \textbf{Erro relativo dos parâmetros:}
    \[
    E_{\theta}
    =
    \frac{
        \left\|
        \hat{\boldsymbol{\theta}}
        -
        \boldsymbol{\theta}^{*}
        \right\|_2
    }{
        \left\|\boldsymbol{\theta}^{*}\right\|_2
    }.
    \]

    \item \textbf{Erro das trajetórias reconstruídas:}
    \[
    \mathrm{RMSE}
    =
    \sqrt{
        \frac{1}{2m}
        \sum_{i=1}^{m}
        \left[
        \left(\hat{N}(t_i)-N^{*}(t_i)\right)^2
        +
        \left(\hat{P}(t_i)-P^{*}(t_i)\right)^2
        \right]
    }.
    \]

    \item \textbf{Estabilidade e robustez:}
    \begin{itemize}
        \item Média e desvio-padrão dos erros
        em múltiplas execuções;
        \item Variação do desempenho em função
        do ruído e da quantidade de observações.
    \end{itemize}

    \item \textbf{Custo computacional:}
    \begin{itemize}
        \item Tempo de treinamento e calibração;
        \item Número de avaliações da função objetivo;
        \item Tempo de inferência, quando aplicável.
    \end{itemize}

\end{itemize}

\end{frame}

\section{Resultados Esperados e Desafios}

\begin{frame}{Resultados Esperados e Limitações}

\begin{block}{Resultados Esperados}
\begin{itemize}

    \item Implementação e avaliação de diferentes
    técnicas de identificação de parâmetros
    para o modelo de Rosenzweig-MacArthur;

    \item Análise comparativa da
    \textbf{precisão, robustez e eficiência}
    dos métodos investigados;

    \item Identificação das vantagens e limitações
    do algoritmo genético e das abordagens baseadas
    em redes neurais;
\end{itemize}
\end{block}

\vspace{0.2cm}

\begin{alertblock}{Limitações e Desafios}
\begin{itemize}

    \item Possíveis dificuldades de identificabilidade
    e sensibilidade dos parâmetros;

    \item Dependência das condições iniciais
    e das configurações de otimização;

    \item Custo de treinamento das redes neurais
    em comparação com a resolução numérica das EDOs.
\end{itemize}
\end{alertblock}

\end{frame}
